% Generated by scripts/assemble.py; corrections belong in source/corrections.json.
\examhead{89}{苏州大学 2025 年数学分析试卷}{88}
% p089-b01; source PDF page 89
\par\Needspace{6\baselineskip}
\blockanchor{p089-b01}
\question{1}设\par
\nopagebreak[4]
% p089-b02; source PDF page 89
\blockanchor{p089-b02}
\displayformula{f (x) = \left\{ \begin{array}{l l} \frac {\sin (x ^ {2})}{x ^ {2}}, & x \neq 0; \\ 1, & x = 0. \end{array} \right.}
% p089-b03; source PDF page 89
\blockanchor{p089-b03}
求 $f^{(2024)}(0)$ .\par
% p089-b04; source PDF page 89
\blockanchor{p089-b04}
\question{2}已知 $f(x) = \cos (x^p)$ 在 $(0, +\infty)$ 上一致连续，求 $p$ 的范围.\par
% p089-b05; source PDF page 89
\blockanchor{p089-b05}
\question{3}设 $I(x) = \int_{0}^{+\infty} e^{-xy^2} \sin x \, \mathrm{d}y, x \geq 0$, 求 $\int_{0}^{+\infty} I(x) \, \mathrm{d}x$.\par
% p089-b06; source PDF page 89
\par\Needspace{6\baselineskip}
\blockanchor{p089-b06}
\question{4}设平面 $x + y + z = 1$ 在球 $x^{2} + y^{2} + z^{2} \leq 1$ 内部的平面部分为 $S$, 计算\par
\nopagebreak[4]
% p089-b07; source PDF page 89
\blockanchor{p089-b07}
\displayformula{\iint_ {S} \left(1 - x ^ {2} - y ^ {2} - z ^ {2}\right) \mathrm{d} S.}
% p089-b08; source PDF page 89
\blockanchor{p089-b08}
\question{5}设 $f(x), g(x)$ 在 $[a, b]$ 上黎曼可积, 证明: $f(x)g(x)$ 在 $[a, b]$ 上黎曼可积.\par
% p089-b09; source PDF page 89
\par\Needspace{6\baselineskip}
\blockanchor{p089-b09}
\question{6}设 $f(x)$ 在 $[0,1]$ 上连续可微, $\int_0^1 f(x)\mathrm{d}x = 0$, 且有\par
\nopagebreak[4]
% p089-b10; source PDF page 89
\blockanchor{p089-b10}
\displayformula{f ^ {\prime} (x) + f (x) \tan x = \int_ {0} ^ {x} f (t) \mathrm{d} t.}
% p089-b11; source PDF page 89
\blockanchor{p089-b11}
求证: $f(x)$ 在 $[0,1]$ 上恒为零.\par
% p089-b12; source PDF page 89
\par\Needspace{6\baselineskip}
\blockanchor{p089-b12}
\question{7}函数 $f(x, y)$ 存在二阶连续偏导数，且 $f(0, 0) = 2024$，$f_{xx} + f_{yy} = 2xy$，记 $L_r: x^2 + y^2 = r^2 (r > 0)$，求\par
\nopagebreak[4]
% p089-b13; source PDF page 89
\blockanchor{p089-b13}
\displayformula{I (r) = \frac {1}{2 \pi r} \int_ {L _ {r}} f (x, y) \mathrm{d} s.}
% p089-b14; source PDF page 89
\par\Needspace{6\baselineskip}
\blockanchor{p089-b14}
\question{8}判断级数\par
\nopagebreak[4]
% p089-b15; source PDF page 89
\blockanchor{p089-b15}
\displayformula{\sum_ {n = 1} ^ {\infty} \frac {(- 1) ^ {n}}{\sqrt {n}}, \sum_ {n = 1} ^ {\infty} \frac {(- 1) ^ {n}}{\sqrt {n} - (- 1) ^ {n}}, \sum_ {n = 1} ^ {\infty} \frac {(- 1) ^ {[ \sqrt {n} ]}}{n}}
% p089-b16; source PDF page 89
\blockanchor{p089-b16}
的敛散性.\par
% p089-b17; source PDF page 89
\blockanchor{p089-b17}
\question{9}设 $f(x)$ 在 $[0,1]$ 上二阶连续可微, $f(0) < 0 < f(1)$, 且 $f'(x) > 0$, $f''(x) > 0$.\par
% p089-b18; source PDF page 89
\blockanchor{p089-b18}
(1) 证明: 存在唯一的 $\xi \in (0,1)$, 满足 $f(\xi) = 0$.\par
% p089-b19; source PDF page 89
\blockanchor{p089-b19}
(2) 记 $x_{1} = 1$, $x_{n+1} = x_{n} - \frac{f(x_{n})}{f'(x_{n})}$, 证明: $\lim_{n \to \infty} x_{n} = \xi$.\par
% p089-b20; source PDF page 89
\blockanchor{p089-b20}
(3) 求极限 $\lim_{n\to \infty}\frac{x_{n + 1} - \xi}{(x_n - \xi)^2}.$\par
